\section{Shapley value (SHAP의 수학)}

2026-10-01. 설명 가능한 AI(XAI)의 SHAP 기법 밑바탕인 협력게임 이론.

\subsection{정의}

특성 집합 \(N\) (\(|N|=n\)), 부분집합 \(S\)만 알 때의 예측값을 가치
함수로 둔다. 빠진 특성은 배경 데이터 \(x'\)로 채워 평균낸다.

\[
v(S)=\mathbb E_{x'}\big[f(x_S,\ x'_{\bar S})\big],\qquad v(\varnothing)=\mathbb E_{x'}[f(x')]
\]

특성 \(i\)의 기여도:

\[
\phi_i=\sum_{S\subseteq N\setminus\{i\}}\frac{|S|!\,(n-|S|-1)!}{n!}\Big[v(S\cup\{i\})-v(S)\Big]
\]

가중치는 ``특성을 무작위 순서로 하나씩 넣을 때 \(i\) 바로 앞에 정확히
\(S\)가 들어와 있을 확률''이다.

효율성 (기여도의 합이 보존됨):

\[
\sum_{i=1}^{n}\phi_i=v(N)-v(\varnothing)
\]

\subsection{2개 특성 예 (집값: 면적 A, 역세권 B)}

\(v(\varnothing)=100,\ v(\{A\})=150,\ v(\{B\})=120,\ v(\{A,B\})=200\)

\[
\phi_A=\tfrac12(150-100)+\tfrac12(200-120)=65,\qquad \phi_B=\tfrac12(120-100)+\tfrac12(200-150)=35
\]

합 \(100=v(\{A,B\})-v(\varnothing)\).

\subsection{\texorpdfstring{연습문제 (\(n=3\))}{연습문제 (n=3)}}

\(v(\varnothing)=10,\ v(A)=20,\ v(B)=10,\ v(C)=10,\ v(AB)=30,\ v(AC)=20,\ v(BC)=10,\ v(ABC)=40\)

\(n=3\)의 가중치: \(|S|=0\)이면 \(\tfrac13\), \(|S|=1\)이면
\(\tfrac16\), \(|S|=2\)이면 \(\tfrac13\).

\[
\phi_A=\tfrac13(10)+\tfrac16(20)+\tfrac16(10)+\tfrac13(30)=\tfrac{55}{3},\qquad \phi_B=\tfrac{25}{3},\qquad \phi_C=\tfrac{10}{3}
\]

합 \(30=v(ABC)-v(\varnothing)\).

\textbf{이상한 점.} \(C\)는 혼자서도, \(A\)나 \(B\) 하나와 함께여도 값을
전혀 바꾸지 않는다. 그런데 \(A,B\)가 이미 있을 때만 \(+10\)을 만든다. 이
3자 상호작용이 셋에게 나눠지면서 \(\phi_C=\tfrac{10}{3}>0\)이 된다.
Shapley value는 상호작용을 쪼개서 개별 기여처럼 보이게 만든다. 이게
SHAP을 ``특성 중요도''로 읽을 때 생기는 비판의 출발점이다.

\subsection{실제 계산}

\(2^n\)개 부분집합을 다 볼 수 없어서 근사한다: KernelSHAP(가중
선형회귀), TreeSHAP(트리 구조로 다항 시간), DeepSHAP(DeepLIFT 기반).
